% workbook.tex — full working record (research journal in LaTeX)
%
% USAGE BY CLAUDE:
%   - This is the COMPLETE working record: proofs, derivations, failed attempts,
%     discussions, corrections, and thinking in progress. Not a paper draft.
%   - Show every step of every derivation. Never summarise with "one finds" or
%     "a short computation gives."
%   - Use \label on every section, subsection, theorem, equation, and figure.
%   - Corrections: REPLACE wrong content in place — never append a correction block.
%   - Grep and targeted reads only; do not read this file top-to-bottom each session.
%
\documentclass[11pt, a4paper]{article}

% --- Core math packages ---
\usepackage{amsmath, amssymb, amsthm}
\usepackage{mathtools}          % \coloneqq, \DeclarePairedDelimiter, etc.
\usepackage{bm}                 % \bm for bold math symbols

% --- Layout ---
\usepackage[margin=2.5cm]{geometry}
\usepackage{microtype}
\usepackage{setspace}
\onehalfspacing

% --- Fonts and encoding ---
\usepackage[T1]{fontenc}
\usepackage[utf8]{inputenc}

% --- References and links ---
\usepackage[colorlinks=true, linkcolor=blue, citecolor=blue, urlcolor=blue]{hyperref}
\usepackage[numbers, sort&compress]{natbib}

% --- Miscellaneous ---
\usepackage{enumitem}
\usepackage{graphicx}
\usepackage{booktabs}           % professional tables

% -----------------------------------------------------------------------
% THEOREM ENVIRONMENTS
% -----------------------------------------------------------------------
\theoremstyle{plain}
\newtheorem{theorem}{Theorem}[section]
\newtheorem{proposition}[theorem]{Proposition}
\newtheorem{lemma}[theorem]{Lemma}
\newtheorem{corollary}[theorem]{Corollary}
\newtheorem{conjecture}[theorem]{Conjecture}

\theoremstyle{definition}
\newtheorem{definition}[theorem]{Definition}
\newtheorem{example}[theorem]{Example}
\newtheorem{remark}[theorem]{Remark}

% -----------------------------------------------------------------------
% CUSTOM COMMANDS
% (add your notation here; keep in sync with brief.tex and CLAUDE.md)
% -----------------------------------------------------------------------
% \newcommand{\examplecommand}[1]{\mathcal{#1}}

% -----------------------------------------------------------------------
\title{[Project Title]}
\author{[Author Name]}
\date{\today}
% -----------------------------------------------------------------------

\begin{document}

\maketitle

\begin{abstract}
[Abstract — one paragraph summary of the main results.]
\end{abstract}

\tableofcontents
\newpage

% -----------------------------------------------------------------------
\section{Introduction}
\label{sec:intro}
% -----------------------------------------------------------------------

[Introduction — background, motivation, statement of main results.]

% -----------------------------------------------------------------------
\section{Notation and conventions}
\label{sec:conventions}
% -----------------------------------------------------------------------

[Define every symbol used in the paper. This section is the reference for
Claude. Be explicit about normalisations, sign conventions, and branch choices.]

% -----------------------------------------------------------------------
\section{[First main topic]}
\label{sec:topic1}
% -----------------------------------------------------------------------

[Content.]

% -----------------------------------------------------------------------
\section{[Second main topic]}
\label{sec:topic2}
% -----------------------------------------------------------------------

[Content.]

% -----------------------------------------------------------------------
\section{Numerical results}
\label{sec:numerics}
% -----------------------------------------------------------------------

[Numerical computations: method, precision, validation, results.]

% -----------------------------------------------------------------------
\section{Conclusion}
\label{sec:conclusion}
% -----------------------------------------------------------------------

[Summary and outlook.]

% -----------------------------------------------------------------------
\bibliographystyle{plainnat}
\bibliography{references}
% -----------------------------------------------------------------------

\end{document}
